\documentclass[11pt,a4paper]{article}
\usepackage{amsmath,amsthm,amssymb,mathtools}
\usepackage[margin=1in]{geometry}

\newtheorem{theorem}{Theorem}
\newtheorem{lemma}[theorem]{Lemma}
\newtheorem{corollary}[theorem]{Corollary}
\newtheorem{definition}[theorem]{Definition}
\newtheorem{remark}[theorem]{Remark}
\newtheorem{proposition}[theorem]{Proposition}

\title{Uniformly Modulated Representation of the Hardy $Z$-Function:\\
Derivation and Recovery of $e^{i\vartheta(t)}\zeta(\tfrac{1}{2}+it)$}
\author{Stephen Crowley}
\date{March 2026}

\begin{document}
\maketitle

\section{Generalized Harmonic Analysis and Spectral Representations of Deterministic Functions}

The spectral representation machinery used throughout this paper does not
depend on probabilistic randomness. This section establishes the precise
functional-analytic framework.

\begin{definition}[Wiener's time-average inner product]
\label{def:wiener-ip}
For locally square-integrable functions $f,g\in L^2_{\mathrm{loc}}(\mathbb{R})$,
define the \emph{time-average inner product}
\begin{equation}
  \langle f,g\rangle_W
  =\lim_{T\to\infty}\frac{1}{2T}\int_{-T}^{T}f(t)\overline{g(t)}\,dt,
\end{equation}
whenever the limit exists. Write
$\|f\|_W^2=\langle f,f\rangle_W$
and $\mathcal{W}=\{f\in L^2_{\mathrm{loc}}(\mathbb{R}):\|f\|_W<\infty\}$.
\end{definition}

\begin{theorem}[Wiener, 1930; Cram\'er--Lo\`eve spectral decomposition]
\label{thm:cramer}
Every $f\in\mathcal{W}$ admits a spectral representation
\begin{equation}
  \label{eq:cramer-rep}
  f(t)=\int_{-\infty}^{\infty}e^{i\omega t}\,d\mu_f(\omega),
\end{equation}
where $d\mu_f$ is the \emph{spectral measure of $f$}, a complex
Borel--Stieltjes measure uniquely determined by
\begin{equation}
  \mu_f((\omega_1,\omega_2])
  =\lim_{T\to\infty}\frac{1}{2T}\int_{-T}^{T}
   f(t)\,\frac{e^{-i\omega_1 t}-e^{-i\omega_2 t}}{it}\,dt.
\end{equation}
The spectral measure has orthogonal increments in the time-average sense:
for disjoint intervals $(\omega_1,\omega_2]$ and $(\omega_3,\omega_4]$,
\begin{equation}
  \lim_{T\to\infty}\frac{1}{2T}\int_{-T}^{T}
  \bigl[e^{-i\omega_1 t}-e^{-i\omega_2 t}\bigr]
  \overline{\bigl[e^{-i\omega_3 t}-e^{-i\omega_4 t}\bigr]}\,dt
  =0.
\end{equation}
The autocorrelation function
\begin{equation}
  R_f(\tau)=\lim_{T\to\infty}\frac{1}{2T}\int_{-T}^{T}
            f(t+\tau)\overline{f(t)}\,dt
\end{equation}
satisfies the Wiener--Khintchine relation
\begin{equation}
  \label{eq:wk}
  R_f(\tau)=\int_{-\infty}^{\infty}e^{i\omega\tau}\,dF_f(\omega),
  \qquad
  dF_f(\omega)=|d\mu_f(\omega)|^2.
\end{equation}
\end{theorem}

\begin{remark}[The notation $\mathbb{E}$ and $dW$]
\label{rem:notation}
For a deterministic function $f$, the symbols
$\mathbb{E}[\cdot]$ and $dW(\omega)$ are notational devices:
\begin{enumerate}
  \item $\mathbb{E}[\cdot]$ denotes the time average
        $\lim_{T\to\infty}\frac{1}{2T}\int_{-T}^{T}(\cdot)\,dt$,
        not a probabilistic expectation.
  \item $dW(\omega)$ denotes an orthogonal-increment measure with
        $\langle dW(\omega),dW(\lambda)\rangle_W=\delta(\omega-\lambda)\,d\omega$,
        encoding the orthogonality of the Fourier basis
        $\{e^{i\omega t}\}_{\omega\in\mathbb{R}}$
        in the time-average inner product.
  \item The Wiener isometry
        $\langle\int g\,dW,\int h\,dW\rangle_W
         =\int g(\omega)\overline{h(\omega)}\,d\omega$
        is the Parseval--Plancherel identity,
        an isometric fact, not a probabilistic one.
\end{enumerate}
When a spectral density $K(\omega)\ge 0$ exists such that
$dF_f(\omega)=K(\omega)\,d\omega$, the spectral factor
$\hat{h}(\omega)=\sqrt{K(\omega)}$ and the colored measure
$d\widetilde{Z}(\omega)=\hat{h}(\omega)\,dW(\omega)$
are defined so that
$|d\widetilde{Z}(\omega)|^2=K(\omega)\,d\omega$
and $f(t)=\int e^{i\omega t}\,d\widetilde{Z}(\omega)$.
This is the standard spectral factorization of generalized harmonic analysis.
\end{remark}

\begin{remark}[Applicability to the Riemann zeta function]
\label{rem:zeta-applicability}
The Hardy $Z$-function
$Z(t)=e^{i\vartheta(t)}\zeta(\tfrac{1}{2}+it)$
is a specific real-valued analytic function of~$t$. Its mean-square
properties over long intervals are described by the second and fourth
moment theorems of Hardy--Littlewood--Ingham--Heath-Brown, and the
Fourier structure of $\zeta(\tfrac{1}{2}+it)$ over the Dirichlet series
provides a natural spectral decomposition. The Cram\'er--Lo\`eve framework
of Theorem~\ref{thm:cramer} applies to $Z$ directly: $Z\in\mathcal{W}$
and all spectral identities that follow are isometric equalities in the
time-average sense. No probability space is required; no randomness is
assumed or introduced.
\end{remark}


\section{Setup and Definitions}

\begin{definition}[Rao spectral density and corrected Bochner identification]
\label{def:rao}
The \emph{Rao spectral density} is
\begin{equation}
  K_{1/2}(\omega)=\frac{e^{\omega/2}}{e^{e^\omega}+1}.
\end{equation}
Starting from the Mellin--Euler integral representation valid for
$\operatorname{Re}s>0$,
\begin{equation}
  \Gamma(s)\,\eta(s)=\int_0^{\infty}\frac{x^{s-1}}{e^x+1}\,dx,
\end{equation}
and applying the change of variables $x=e^{\omega}$,
$dx=e^{\omega}\,d\omega$, $x^{s-1}=e^{\omega(s-1)}$, one obtains at
$s=\tfrac{1}{2}+i\tau$:
\begin{equation}
  \label{eq:bochner}
  \Gamma\!\left(\tfrac{1}{2}+i\tau\right)
  \eta\!\left(\tfrac{1}{2}+i\tau\right)
  =\int_{-\infty}^{\infty} e^{i\omega\tau}\,K_{1/2}(\omega)\,d\omega.
\end{equation}
Since $K_{1/2}\ge 0$ and $K_{1/2}\in L^1(\mathbb{R})$, Bochner's theorem
identifies $K_{1/2}$ as the spectral density of the positive-definite
kernel
$\tau\mapsto\Gamma(\tfrac{1}{2}+i\tau)\,\eta(\tfrac{1}{2}+i\tau)$.
The product $\Gamma(\tfrac{1}{2}+i\tau)\eta(\tfrac{1}{2}+i\tau)$
decays as $e^{-\pi|\tau|/2}$ by Stirling's formula, so it is integrable and
the Fourier inversion is unconditional.
\end{definition}

\begin{definition}[Spectral factor and orthogonal measure]
\label{def:spectral-factor}
Set
\begin{equation}
  \hat{h}(\omega)=\sqrt{K_{1/2}(\omega)},
  \qquad
  d\widetilde{Z}(\omega)=\hat{h}(\omega)\,dW(\omega),
\end{equation}
where $dW(\omega)$ is an orthogonal-increment measure with
$|dW(\omega)|^2=d\omega$ (see Remark~\ref{rem:notation}), so that
$|d\widetilde{Z}(\omega)|^2=K_{1/2}(\omega)\,d\omega$.
\end{definition}

\begin{definition}[Stationary core]
\label{def:stationary-core}
Define the \emph{stationary core} by
\begin{equation}
  \label{eq:S-def}
  S(\tau)=\int_{-\infty}^{\infty} e^{i\omega\tau}\,d\widetilde{Z}(\omega)
  =\int_{-\infty}^{\infty} e^{i\omega\tau}\hat{h}(\omega)\,dW(\omega).
\end{equation}
By the Wiener isometry (Parseval identity) and \eqref{eq:bochner},
\begin{equation}
  \label{eq:S-covariance}
  \langle S(\tau),S(\sigma)\rangle_W
  =\int_{-\infty}^{\infty} e^{i\omega(\tau-\sigma)}K_{1/2}(\omega)\,d\omega
  =\Gamma\!\left(\tfrac{1}{2}+i(\tau-\sigma)\right)
   \eta\!\left(\tfrac{1}{2}+i(\tau-\sigma)\right).
\end{equation}
\end{definition}

\begin{definition}[Riemann--Siegel quantities]
\label{def:rs}
The Riemann--Siegel theta function is
\begin{equation}
  \vartheta(t)=\operatorname{Im}\log\Gamma\!\left(\tfrac{1}{4}+\tfrac{it}{2}\right)
  -\frac{t}{2}\log\pi,
  \qquad
  \vartheta'(t)=\frac{1}{2}\log\frac{t}{2\pi}+O(t^{-2}).
\end{equation}
We treat $\vartheta$ as strictly increasing (valid for $t$ sufficiently
large) and write $\vartheta^{-1}$ for its functional inverse.
The Hardy $Z$-function is
\begin{equation}
  Z(t)=e^{i\vartheta(t)}\zeta\!\left(\tfrac{1}{2}+it\right).
\end{equation}
\end{definition}


\section{Oscillatory Process Structure}

\begin{definition}[Oscillatory kernel and frequency-dependent gain]
\label{def:oscillatory}
In the spectral integral
\begin{equation}
  \label{eq:osc-full}
  Z(t)=\int_{-\infty}^{\infty}
       \phi_t(\omega)\,d\widetilde{Z}(\omega),
\end{equation}
define the \emph{oscillatory kernel}
\begin{equation}
  \label{eq:phi-def}
  \phi_t(\omega)
  =\sqrt{\vartheta'(t)}\,e^{i\omega\vartheta(t)}
  =A_t(\omega)\,e^{i\omega t},
\end{equation}
where the \emph{gain function} is
\begin{equation}
  \label{eq:gain}
  A_t(\omega)=\sqrt{\vartheta'(t)}\,e^{i\omega(\vartheta(t)-t)}.
\end{equation}
Since $\vartheta(t)-t=\tfrac{t}{2}\log\tfrac{t}{2\pi}-\tfrac{t}{2}
-\tfrac{\pi}{8}+O(t^{-1})-t\neq 0$ for all finite~$t$, the gain
$A_t(\omega)$ depends on $\omega$ through the phase factor
$e^{i\omega(\vartheta(t)-t)}$. This makes \eqref{eq:osc-full} an
\emph{oscillatory process} in the sense of Priestley (1965), not merely
a uniformly modulated one.
\end{definition}

\begin{remark}[Why the distinction matters]
A \emph{uniformly modulated} process has $A_t(\omega)=c(t)$ independent
of~$\omega$, so modulation is pure amplitude scaling. Here, the
$\omega$-dependent phase $e^{i\omega(\vartheta(t)-t)}$ performs a
\emph{local frequency remapping}: each spectral component at
frequency~$\omega$ is shifted in phase by
$\omega\cdot(\vartheta(t)-t)$, which grows with $\omega$.
The amplitude factor $\sqrt{\vartheta'(t)}$ provides the Jacobian
normalization for the time change $t\mapsto\vartheta(t)$, while the
phase factor $e^{i\omega(\vartheta(t)-t)}$ encodes the nonlinear
warping of the frequency axis. Together they implement the full
pullback from the $\tau$-coordinate to the $t$-coordinate.
\end{remark}


\section{The Representation and Recovery Theorem}

\begin{theorem}
\label{thm:main}
With the stationary core $S$ of Definition~\ref{def:stationary-core},
\begin{equation}
  \label{eq:Zrep}
  Z(t)=\sqrt{\vartheta'(t)}\,S(\vartheta(t)).
\end{equation}
Expanding $S$ via its spectral representation \eqref{eq:S-def},
\begin{equation}
  \label{eq:Zfull}
  Z(t)=\sqrt{\vartheta'(t)}\int_{-\infty}^{\infty}
        e^{i\omega\vartheta(t)}\hat{h}(\omega)\,dW(\omega)
  =\int_{-\infty}^{\infty}A_t(\omega)\,e^{i\omega t}\,
   \hat{h}(\omega)\,dW(\omega).
\end{equation}
This equals $e^{i\vartheta(t)}\zeta(\tfrac{1}{2}+it)$ as a
pointwise identity in $t$, where both sides are specific analytic
functions and all spectral integrals are Fourier--Stieltjes integrals
in the sense of Theorem~\ref{thm:cramer}.
\end{theorem}

\begin{proof}
\medskip\noindent\textbf{Step 1: evaluate $S$ at $\tau=\vartheta(t)$.}
Substituting $\tau=\vartheta(t)$ into \eqref{eq:S-def},
\begin{equation}
  \label{eq:S-at-theta}
  S(\vartheta(t))
  =\int_{-\infty}^{\infty} e^{i\omega\vartheta(t)}\,d\widetilde{Z}(\omega).
\end{equation}

\medskip\noindent\textbf{Step 2: form the oscillatory representation.}
Multiply by $\sqrt{\vartheta'(t)}$:
\begin{equation}
  \label{eq:Z-expand}
  Z(t)
  =\sqrt{\vartheta'(t)}\,S(\vartheta(t))
  =\sqrt{\vartheta'(t)}\int_{-\infty}^{\infty}
   e^{i\omega\vartheta(t)}\hat{h}(\omega)\,dW(\omega)
  =\int_{-\infty}^{\infty}
   A_t(\omega)\,e^{i\omega t}\,d\widetilde{Z}(\omega).
\end{equation}

\medskip\noindent\textbf{Step 3: compute the autocorrelation.}
Apply the Wiener isometry (Parseval identity):
\begin{align}
  \langle Z(t+\cdot),Z(s+\cdot)\rangle_W
  &=\sqrt{\vartheta'(t)\vartheta'(s)}
    \int_{-\infty}^{\infty}
    e^{i\omega(\vartheta(t)-\vartheta(s))}K_{1/2}(\omega)\,d\omega \\
  \label{eq:Zcov}
  &=\sqrt{\vartheta'(t)\vartheta'(s)}\,
   \Gamma\!\left(\tfrac{1}{2}+i(\vartheta(t)-\vartheta(s))\right)
   \eta\!\left(\tfrac{1}{2}+i(\vartheta(t)-\vartheta(s))\right).
\end{align}

\medskip\noindent\textbf{Step 4: connection to $\zeta$ via $\eta$.}
The Dirichlet eta function and Riemann zeta function satisfy
\begin{equation}
  \label{eq:eta-zeta}
  \eta(s)=(1-2^{1-s})\zeta(s).
\end{equation}
At $s=\tfrac{1}{2}+it$, the Bochner identification \eqref{eq:bochner}
gives
\begin{equation}
  \Gamma\!\left(\tfrac{1}{2}+it\right)
  \eta\!\left(\tfrac{1}{2}+it\right)
  =\int_{-\infty}^{\infty} e^{i\omega t}\,K_{1/2}(\omega)\,d\omega,
\end{equation}
hence
\begin{equation}
  \eta\!\left(\tfrac{1}{2}+it\right)
  =\frac{1}{\Gamma(\tfrac{1}{2}+it)}
   \int_{-\infty}^{\infty} e^{i\omega t}\,K_{1/2}(\omega)\,d\omega
\end{equation}
and
\begin{equation}
  \zeta\!\left(\tfrac{1}{2}+it\right)
  =\frac{1}{(1-2^{1/2-it})\,\Gamma(\tfrac{1}{2}+it)}
   \int_{-\infty}^{\infty} e^{i\omega t}\,K_{1/2}(\omega)\,d\omega.
\end{equation}
Multiplying by $e^{i\vartheta(t)}$:
\begin{equation}
  \label{eq:Z-deterministic}
  Z(t)=e^{i\vartheta(t)}\zeta\!\left(\tfrac{1}{2}+it\right)
  =\frac{e^{i\vartheta(t)}}{(1-2^{1/2-it})\,\Gamma(\tfrac{1}{2}+it)}
   \int_{-\infty}^{\infty} e^{i\omega t}\,K_{1/2}(\omega)\,d\omega.
\end{equation}

\medskip\noindent\textbf{Step 5: the bridge.}
Both \eqref{eq:Zfull} and \eqref{eq:Z-deterministic} express $Z(t)$
as Fourier--Stieltjes integrals built from the same spectral density
$K_{1/2}$. The oscillatory representation \eqref{eq:Zfull} uses the
warped phase $e^{i\omega\vartheta(t)}$ and the Jacobian prefactor
$\sqrt{\vartheta'(t)}$; the deterministic representation
\eqref{eq:Z-deterministic} uses the standard phase $e^{i\omega t}$ and
the algebraic prefactor
$e^{i\vartheta(t)}/[(1-2^{1/2-it})\Gamma(\tfrac{1}{2}+it)]$.
They are the same function evaluated in two coordinate systems related
by the diffeomorphism $t\mapsto\vartheta(t)$.

Equating them gives the explicit identity for $S(\vartheta(t))$:
\begin{equation}
  \label{eq:bridge}
  S(\vartheta(t))
  =\frac{e^{i\vartheta(t)}}
        {\sqrt{\vartheta'(t)}\,(1-2^{1/2-it})\,\Gamma(\tfrac{1}{2}+it)}
   \int_{-\infty}^{\infty} e^{i\omega t}\,K_{1/2}(\omega)\,d\omega
  =\frac{e^{i\vartheta(t)}}{\sqrt{\vartheta'(t)}}\,
   \zeta\!\left(\tfrac{1}{2}+it\right),
\end{equation}
so that
\begin{equation}
  Z(t)=\sqrt{\vartheta'(t)}\,S(\vartheta(t))
      =e^{i\vartheta(t)}\zeta\!\left(\tfrac{1}{2}+it\right).
\end{equation}
\end{proof}


\section{Recapitulation: The Four Layers}

\noindent\textbf{1.\ Orthogonal-increment measure:}
\[
  |dW(\omega)|^2=d\omega.
\]

\noindent\textbf{2.\ Spectral coloring:}
\[
  d\widetilde{Z}(\omega)=\hat{h}(\omega)\,dW(\omega),
  \qquad
  |d\widetilde{Z}(\omega)|^2=K_{1/2}(\omega)\,d\omega.
\]

\noindent\textbf{3.\ Stationary core in the $\tau$-coordinate:}
\[
  S(\tau)=\int_{-\infty}^{\infty} e^{i\omega\tau}\,d\widetilde{Z}(\omega),
  \qquad
  \langle S(\tau),S(\sigma)\rangle_W
  =\Gamma\!\left(\tfrac{1}{2}+i(\tau-\sigma)\right)
   \eta\!\left(\tfrac{1}{2}+i(\tau-\sigma)\right).
\]

\noindent\textbf{4.\ Oscillatory pullback to the $t$-coordinate:}
\[
  Z(t)=\sqrt{\vartheta'(t)}\,S(\vartheta(t))
  =\int_{-\infty}^{\infty}A_t(\omega)\,e^{i\omega t}\,d\widetilde{Z}(\omega),
  \qquad
  A_t(\omega)=\sqrt{\vartheta'(t)}\,e^{i\omega(\vartheta(t)-t)}.
\]

\noindent\textbf{Pointwise recovery:}
\[
  Z(t)=e^{i\vartheta(t)}\zeta\!\left(\tfrac{1}{2}+it\right)
  =\frac{e^{i\vartheta(t)}}{(1-2^{1/2-it})\,\Gamma(\tfrac{1}{2}+it)}
   \int_{-\infty}^{\infty} e^{i\omega t}\,K_{1/2}(\omega)\,d\omega.
\]

\noindent\textbf{Autocorrelation:}
\[
  \langle Z(t+\cdot),Z(s+\cdot)\rangle_W
  =\sqrt{\vartheta'(t)\vartheta'(s)}\,
   \Gamma\!\left(\tfrac{1}{2}+i(\vartheta(t)-\vartheta(s))\right)
   \eta\!\left(\tfrac{1}{2}+i(\vartheta(t)-\vartheta(s))\right).
\]

\end{document}
 